%Version 3 October 2023
%DIF LATEXDIFF DIFFERENCE FILE
%DIF DEL manuscript/article_overleaf.tex   Mon Jun 17 16:59:12 2024
%DIF ADD manuscript/sn-article.tex         Wed Jun 19 14:08:35 2024
% See section 11 of the User Manual for version history
%
%%%%%%%%%%%%%%%%%%%%%%%%%%%%%%%%%%%%%%%%%%%%%%%%%%%%%%%%%%%%%%%%%%%%%%
%%                                                                 %%
%% Please do not use \input{...} to include other tex files.       %%
%% Submit your LaTeX manuscript as one .tex document.              %%
%%                                                                 %%
%% All additional figures and files should be attached             %%
%% separately and not embedded in the \TeX\ document itself.       %%
%%                                                                 %%
%%%%%%%%%%%%%%%%%%%%%%%%%%%%%%%%%%%%%%%%%%%%%%%%%%%%%%%%%%%%%%%%%%%%%

%%\documentclass[referee,sn-basic]{sn-jnl}% referee option is meant for double line spacing

%%=======================================================%%
%% to print line numbers in the margin use lineno option %%
%%=======================================================%%

\documentclass[lineno,sn-basic]{sn-jnl}% Basic Springer Nature Reference Style/Chemistry Reference Style

%%======================================================%%
%% to compile with pdflatex/xelatex use pdflatex option %%
%%======================================================%%

%%\documentclass[pdflatex,sn-basic]{sn-jnl}% Basic Springer Nature Reference Style/Chemistry Reference Style


%%Note: the following reference styles support Namedate and Numbered referencing. By default the style follows the most common style. To switch between the options you can add or remove “Numbered” in the optional parenthesis. 
%%The option is available for: sn-basic.bst, sn-vancouver.bst, sn-chicago.bst%  
 
%%\documentclass[sn-nature]{sn-jnl}% Style for submissions to Nature Portfolio journals
%%\documentclass[sn-basic]{sn-jnl}% Basic Springer Nature Reference Style/Chemistry Reference Style
% \documentclass[sn-mathphys-num]{sn-jnl}% Math and Physical Sciences Numbered Reference Style 
%%\documentclass[sn-mathphys-ay]{sn-jnl}% Math and Physical Sciences Author Year Reference Style
%%\documentclass[sn-aps]{sn-jnl}% American Physical Society (APS) Reference Style
%%\documentclass[sn-vancouver,Numbered]{sn-jnl}% Vancouver Reference Style
%%\documentclass[sn-apa]{sn-jnl}% APA Reference Style 
%%\documentclass[sn-chicago]{sn-jnl}% Chicago-based Humanities Reference Style

%%%% Standard Packages
%%<additional latex packages if required can be included here>

\usepackage{graphicx}%
\usepackage{multirow}%
\usepackage{amsmath,amssymb,amsfonts}%
\usepackage{amsthm}%
\usepackage{mathrsfs}%
\usepackage[title]{appendix}%
\usepackage{xcolor}%
\usepackage{textcomp}%
\usepackage{manyfoot}%
\usepackage{booktabs}%
\usepackage{algorithm}%
\usepackage{algorithmicx}%
\usepackage{algpseudocode}%
\usepackage{listings}%
% \usepackage{natbib}%

\raggedbottom
%%\unnumbered% uncomment this for unnumbered level heads
%DIF PREAMBLE EXTENSION ADDED BY LATEXDIFF
%DIF UNDERLINE PREAMBLE %DIF PREAMBLE
\RequirePackage[normalem]{ulem} %DIF PREAMBLE
\RequirePackage{color}\definecolor{RED}{rgb}{1,0,0}\definecolor{BLUE}{rgb}{0,0,1} %DIF PREAMBLE
\providecommand{\DIFadd}[1]{{\protect\color{blue}\uwave{#1}}} %DIF PREAMBLE
\providecommand{\DIFdel}[1]{{\protect\color{red}\sout{#1}}}                      %DIF PREAMBLE
%DIF SAFE PREAMBLE %DIF PREAMBLE
\providecommand{\DIFaddbegin}{} %DIF PREAMBLE
\providecommand{\DIFaddend}{} %DIF PREAMBLE
\providecommand{\DIFdelbegin}{} %DIF PREAMBLE
\providecommand{\DIFdelend}{} %DIF PREAMBLE
\providecommand{\DIFmodbegin}{} %DIF PREAMBLE
\providecommand{\DIFmodend}{} %DIF PREAMBLE
%DIF FLOATSAFE PREAMBLE %DIF PREAMBLE
\providecommand{\DIFaddFL}[1]{\DIFadd{#1}} %DIF PREAMBLE
\providecommand{\DIFdelFL}[1]{\DIFdel{#1}} %DIF PREAMBLE
\providecommand{\DIFaddbeginFL}{} %DIF PREAMBLE
\providecommand{\DIFaddendFL}{} %DIF PREAMBLE
\providecommand{\DIFdelbeginFL}{} %DIF PREAMBLE
\providecommand{\DIFdelendFL}{} %DIF PREAMBLE
\newcommand{\DIFscaledelfig}{0.5}
%DIF HIGHLIGHTGRAPHICS PREAMBLE %DIF PREAMBLE
\RequirePackage{settobox} %DIF PREAMBLE
\RequirePackage{letltxmacro} %DIF PREAMBLE
\newsavebox{\DIFdelgraphicsbox} %DIF PREAMBLE
\newlength{\DIFdelgraphicswidth} %DIF PREAMBLE
\newlength{\DIFdelgraphicsheight} %DIF PREAMBLE
% store original definition of \includegraphics %DIF PREAMBLE
\LetLtxMacro{\DIFOincludegraphics}{\includegraphics} %DIF PREAMBLE
\newcommand{\DIFaddincludegraphics}[2][]{{\color{blue}\fbox{\DIFOincludegraphics[#1]{#2}}}} %DIF PREAMBLE
\newcommand{\DIFdelincludegraphics}[2][]{% %DIF PREAMBLE
\sbox{\DIFdelgraphicsbox}{\DIFOincludegraphics[#1]{#2}}% %DIF PREAMBLE
\settoboxwidth{\DIFdelgraphicswidth}{\DIFdelgraphicsbox} %DIF PREAMBLE
\settoboxtotalheight{\DIFdelgraphicsheight}{\DIFdelgraphicsbox} %DIF PREAMBLE
\scalebox{\DIFscaledelfig}{% %DIF PREAMBLE
\parbox[b]{\DIFdelgraphicswidth}{\usebox{\DIFdelgraphicsbox}\\[-\baselineskip] \rule{\DIFdelgraphicswidth}{0em}}\llap{\resizebox{\DIFdelgraphicswidth}{\DIFdelgraphicsheight}{% %DIF PREAMBLE
\setlength{\unitlength}{\DIFdelgraphicswidth}% %DIF PREAMBLE
\begin{picture}(1,1)% %DIF PREAMBLE
\thicklines\linethickness{2pt} %DIF PREAMBLE
{\color[rgb]{1,0,0}\put(0,0){\framebox(1,1){}}}% %DIF PREAMBLE
{\color[rgb]{1,0,0}\put(0,0){\line( 1,1){1}}}% %DIF PREAMBLE
{\color[rgb]{1,0,0}\put(0,1){\line(1,-1){1}}}% %DIF PREAMBLE
\end{picture}% %DIF PREAMBLE
}\hspace*{3pt}}} %DIF PREAMBLE
} %DIF PREAMBLE
\LetLtxMacro{\DIFOaddbegin}{\DIFaddbegin} %DIF PREAMBLE
\LetLtxMacro{\DIFOaddend}{\DIFaddend} %DIF PREAMBLE
\LetLtxMacro{\DIFOdelbegin}{\DIFdelbegin} %DIF PREAMBLE
\LetLtxMacro{\DIFOdelend}{\DIFdelend} %DIF PREAMBLE
\DeclareRobustCommand{\DIFaddbegin}{\DIFOaddbegin \let\includegraphics\DIFaddincludegraphics} %DIF PREAMBLE
\DeclareRobustCommand{\DIFaddend}{\DIFOaddend \let\includegraphics\DIFOincludegraphics} %DIF PREAMBLE
\DeclareRobustCommand{\DIFdelbegin}{\DIFOdelbegin \let\includegraphics\DIFdelincludegraphics} %DIF PREAMBLE
\DeclareRobustCommand{\DIFdelend}{\DIFOaddend \let\includegraphics\DIFOincludegraphics} %DIF PREAMBLE
\LetLtxMacro{\DIFOaddbeginFL}{\DIFaddbeginFL} %DIF PREAMBLE
\LetLtxMacro{\DIFOaddendFL}{\DIFaddendFL} %DIF PREAMBLE
\LetLtxMacro{\DIFOdelbeginFL}{\DIFdelbeginFL} %DIF PREAMBLE
\LetLtxMacro{\DIFOdelendFL}{\DIFdelendFL} %DIF PREAMBLE
\DeclareRobustCommand{\DIFaddbeginFL}{\DIFOaddbeginFL \let\includegraphics\DIFaddincludegraphics} %DIF PREAMBLE
\DeclareRobustCommand{\DIFaddendFL}{\DIFOaddendFL \let\includegraphics\DIFOincludegraphics} %DIF PREAMBLE
\DeclareRobustCommand{\DIFdelbeginFL}{\DIFOdelbeginFL \let\includegraphics\DIFdelincludegraphics} %DIF PREAMBLE
\DeclareRobustCommand{\DIFdelendFL}{\DIFOaddendFL \let\includegraphics\DIFOincludegraphics} %DIF PREAMBLE
%DIF COLORLISTINGS PREAMBLE %DIF PREAMBLE
\RequirePackage{listings} %DIF PREAMBLE
\RequirePackage{color} %DIF PREAMBLE
\lstdefinelanguage{DIFcode}{ %DIF PREAMBLE
%DIF DIFCODE_UNDERLINE %DIF PREAMBLE
  moredelim=[il][\color{red}\sout]{\%DIF\ <\ }, %DIF PREAMBLE
  moredelim=[il][\color{blue}\uwave]{\%DIF\ >\ } %DIF PREAMBLE
} %DIF PREAMBLE
\lstdefinestyle{DIFverbatimstyle}{ %DIF PREAMBLE
	language=DIFcode, %DIF PREAMBLE
	basicstyle=\ttfamily, %DIF PREAMBLE
	columns=fullflexible, %DIF PREAMBLE
	keepspaces=true %DIF PREAMBLE
} %DIF PREAMBLE
\lstnewenvironment{DIFverbatim}{\lstset{style=DIFverbatimstyle}}{} %DIF PREAMBLE
\lstnewenvironment{DIFverbatim*}{\lstset{style=DIFverbatimstyle,showspaces=true}}{} %DIF PREAMBLE
%DIF END PREAMBLE EXTENSION ADDED BY LATEXDIFF

\begin{document}

% \title[Article Title]{Does the presence of leads in the Arctic sea ice have an impact on the ocean?}

\title[Article Title]{Is the surface forcing through sea ice leads transferred to the Arctic \DIFdelbegin \DIFdel{ocean }\DIFdelend \DIFaddbegin \DIFadd{Ocean }\DIFaddend interior?}
%%=============================================================%%
%% GivenName	-> \fnm{Joergen W.}
%% Particle	-> \spfx{van der} -> surname prefix
%% FamilyName	-> \sur{Ploeg}
%% Suffix	-> \sfx{IV}
%% \author*[1,2]{\fnm{Joergen W.} \spfx{van der} \sur{Ploeg} 
%%  \sfx{IV}}\email{iauthor@gmail.com}
%%=============================================================%%

\author*[1]{\fnm{Josu\'e} \sur{Mart\'inez-Moreno}}\email{josue.martinez.moreno@ifremer.fr}

\author[1]{\fnm{Camille} \sur{Lique}}\email{camille.lique@ifremer.fr}
%\equalcont{These authors contributed equally to this work.}

\author[1]{\fnm{Claude} \sur{Talandier}}\email{claude.Talandier@ifremer.fr}
%\equalcont{These authors contributed equally to this work.}

\affil*[1]{\orgdiv{Univ Brest, CNRS, Ifremer, IRD }, \orgname{Laboratoire d'Oc\'eanographie Physique et Spatiale (LOPS)}, \orgaddress{Institut Universitaire Européen de la Mer, \city{Plouzan\'e}, \postcode{F29280}, \state{Bretagne}, \country{France}}}


%%==================================%%
%% Sample for unstructured abstract %%
%%==================================%%

\DIFdelbegin %DIFDELCMD < \abstract{The Arctic sea-ice, in particular the ice-pack, acts as an insulator between the atmosphere and the ocean below. Leads, common in Arctic sea ice, facilitate exchanges of heat, salt, moisture and gases between the ocean and the atmosphere. 
%DIFDELCMD < Local observations have captured heat fluxes through some leads at least an order of magnitude larger than outside of the leads, leading to the speculation that air-sea exchanges through leads may contribute most of the ocean surface buoyancy forcing.
%DIFDELCMD < %Existing literature suggests that leads play a significant role in the heat budget of the surface of the ocean and likely impact water mass transformation within the Arctic. 
%DIFDELCMD < However, due to the lack of observations and high resolution models capable to resolve leads, an estimate of the impact of leads on the ocean water mass transformation has been elusive.  Here, we quantify the seasonal magnitude and impact on the ocean surface caused by the leads using SEDNA, a subkilometer pan-Arctic hindcast (1/60$^\circ$; NEMO model). Leads account for 22\% of the sea-ice cover extent, and within them, there is approximately 25\% of the total surface water mass transformation. 
%DIFDELCMD < The buoyancy flux and water mass transformation in leads are similar to those underneath the surrounding ice covered oceans and the percentage of water mass transformation through leads matches their area coverage. 
%DIFDELCMD < Thus, the present estimate indicates that leads have a small contribution to Arctic Ocean dynamics, contrary to previous hypotheses.}
%DIFDELCMD < %%%
\DIFdelend \DIFaddbegin \abstract{The Arctic sea ice, in particular the ice pack, acts as an insulator between the atmosphere and the ocean below. Leads, common in Arctic sea ice, facilitate exchanges of heat, salt, moisture, and gases between the ocean and the atmosphere. Local observations have captured heat fluxes through some leads at least an order of magnitude larger than those outside of the leads, leading to the speculation that air-sea exchanges through leads may contribute significantly to the Arctic Ocean surface buoyancy forcing. However, due to the lack of observations and high-resolution models capable of resolving leads, an estimate of the impact of leads on the ocean buoyancy forcing and water mass transformation has been elusive.  Here, we quantify the seasonal magnitude and impact on the ocean surface caused by the leads using SEDNA, a subkilometer pan-Arctic hindcast (1/60$^\circ$; NEMO model). Leads account for 22\% of the sea-ice cover extent, and within them, there is approximately 25\% of the total surface water mass transformation. The buoyancy flux and water mass transformation in leads are similar to those underneath the surrounding ice-covered oceans and the percentage of water mass transformation through leads matches their area coverage. Thus, the present estimate indicates that leads have a small contribution to Arctic Ocean dynamics, contrary to previous hypotheses.}
\DIFaddend 

\keywords{Arctic, Leads, Buoyancy flux}

%%\pacs[JEL Classification]{D8, H51}
%%\pacs[MSC Classification]{35A01, 65L10, 65L12, 65L20, 65L70}

\maketitle

%\section{Introduction}\label{sec1}

The Arctic sea ice regulates Earth's climate by acting as a natural insulator between the atmosphere and the ocean \citep{Wettlaufer_insulator_1997, Untersteiner_ice_insulator_1961}. \DIFdelbegin \DIFdel{An }\DIFdelend \DIFaddbegin \DIFadd{A }\DIFaddend ubiquitous feature of Arctic sea ice is the formation and persistence of leads. Leads occur across the polar oceans, both in the marginal ice zone (MIZ; $15-80\%$ sea ice concentration) and the ice pack ($>80\%$ sea ice concentration). Within the MIZ, leads are primarily formed by the advection of sea ice, while in the ice pack, leads form through deformation of the sea ice and are commonly \DIFdelbegin \DIFdel{refereed }\DIFdelend \DIFaddbegin \DIFadd{referred }\DIFaddend to in literature as ``linear kinematic features" (LKF). 
The primary forcing of the sea ice advection and deformation is the wind \citep{Linow_lkf_2017, Hutter_scaling_lkf_2018}, with a smaller contribution from ocean currents \citep{Willmes_leads_currents_2023}.
Leads result in openings of the sea ice cover \citep{Rampal_nextsim_2016, Hutchings_model_lkf_2005, Richter_pack_stress_2002} with spatial scales of meters to kilometers in width and a few kilometers up to hundreds of kilometers in length, and temporal scales ranging from a few hours up to a few days \citep{Linow_lkf_2017, Wernecke_Lead_2015, Tschudi_obs_leads_1998}. Thus, they can locally impact the ocean surface forcing in the \DIFdelbegin \DIFdel{ice covered }\DIFdelend \DIFaddbegin \DIFadd{ice-covered }\DIFaddend polar oceans \citep{Lupkes_leads_2008}.

Once leads are formed, atmospheric forcing in conjunction with localized upwelling or downwelling occurring in the surface layer can result in important heat fluxes at the ocean surface ($\mathcal{O} \sim 100 W/m^2$; \citealp{Bourgault_flux_lkf_2020,Marcq_flux_lead_2012,McPhee_LKF_2005,Maykut_heat_1986}), instigating localized melting or freezing of sea ice \citep{Albedyll_thickness_dis_2022}. For example, \citet{Boutin_rheology_2023} estimated that 35\% of the total sea ice growth during winter occurs within the leads in the ice pack. The Arctic \DIFdelbegin \DIFdel{ice covered }\DIFdelend \DIFaddbegin \DIFadd{ice-covered }\DIFaddend ocean is a $\beta$-ocean, i.e. the ocean stratification is mainly controlled by salinity, which in turn is largely determined by ice-ocean interactions (sea-ice growth and melt). Thus, leads experiencing an increase in buoyancy forcing due to brine rejection can induce convection \citep{Smith_convection_1998}, weaken the mixed layer stratification, and generate fronts, mixed layer turbulence, and eddies \citep{Reiser_LKF_2020, Smith_generation_2002}. Meanwhile, leads experiencing melting will increase the mixed layer stratification and stabilize the ocean surface layer. These changes in the buoyancy flux translate into local transformation of the surface water masses that could be critical for the functioning of the Arctic Ocean circulation \citep{Lenn_mixing_2022, Pemberton_WMT_2015, Smith_polynyas_1990}. 
Previous studies have shown that instances of leads impact locally the ocean surface and the properties of the Arctic Ocean mixed layer. Nonetheless, the integrated contribution of leads to the \DIFdelbegin \DIFdel{large scale }\DIFdelend \DIFaddbegin \DIFadd{large-scale }\DIFaddend buoyancy forcing at the Arctic \DIFdelbegin \DIFdel{ocean }\DIFdelend \DIFaddbegin \DIFadd{Ocean }\DIFaddend surface has not been quantified yet.

Here, we present the first estimate of the contribution of the buoyancy forcing and water mass transformation within leads in the Arctic Ocean, using the last year of a seven year long pan-Arctic hindcast run at a subkilometer resolution (SEDNA, \citealp{SEDNA_2023}; see Methods for more details). Our main objectives are to: (1) assess the impact of leads in the surface buoyancy forcing and (2) contrast the surface water mass transformation within the leads and those of the ice cover excluding leads.

\section{Buoyancy forcing in leads}\label{sec2}

\begin{figure}
    \centering
    \includegraphics[width=0.95\textwidth]{./Figures/Figure_1_SENTINE_SEDNA_LKF_conc_New_Color.pdf}
    \caption{a) Snapshot of total deformation for the 21st of April 2014 from SEDNA. b) Mean lead probability for daily snapshots of leads identified between 2011 and 2015 from SEDNA. Zoom of the magenta box of panel a for the c) total deformation to identify leads, d) ice concentration, and f) identified mask of leads on the 21st of April 2014 in SEDNA. e) True color image from Sentinel 3 on the 12th of April 2022 for the magenta area in panel a.}
    \label{fig:fig1}
\end{figure}

Leads can be easily visually identified when examining high resolution satellite observations (Figure \ref{fig:fig1}f), yet it remains challenging to capture them in state-of-the-art models that often lack the resolution and/or the sea ice dynamics to simulate these features \citep{Bouchat_sea_ice_rheology_2022, Wang_leads_2016}. Here we analyze daily outputs from the SEDNA ocean-sea ice model, which has an average Arctic resolution of $\sim 800m$, sufficient to capture the observed lead width distribution ($>1km$; \citealp{Wernecke_Lead_2015}). 
We start our analysis by detecting the leads in the model sea ice fields, using the detection algorithm of \citet{Hutter_identification_2019}, that \DIFdelbegin \DIFdel{identify }\DIFdelend \DIFaddbegin \DIFadd{identifies }\DIFaddend leads based on the sea ice deformation field (see methods). Figure \ref{fig:fig1}a and c shows a snapshot of the sea ice deformation for the 21st of April 2014, from which we retrieve a mask of the leads (Figure \ref{fig:fig1}e). Within the leads, the sea ice concentration can remain higher than 90\%; this is a consequence of \DIFaddbegin \DIFadd{the }\DIFaddend continuous nature of the sea ice model and the daily averaged output that smooths the presence of leads in the fields.

The SEDNA climatology of the lead probability, i.e. the likelihood of leads occurring for any given day, is shown in figure \ref{fig:fig1}b. The average over the Arctic amounts to $\sim 19\%$, comparable to previous observations \citep{Wernecke_Lead_2015,Willmes_leads_currents_2023}. Although observations use different identification methods compared to the one implemented in SEDNA \citep{Hutter_identification_2019}, there is resemblance between SEDNA and observations (Fig. \ref{fig:S_fig1}), with similar lead probability over the shelves, along bathymetric features, and in regions with large kinetic energy (Fig. \ref{fig:S_fig1}b). SEDNA reproduces the overall spatial patterns of lead probability observed in the Arctic, except for the Beaufort Sea. Regardless, the similarities between the observations and the model lead probability give us confidence to quantify for the first time the impact of leads \DIFdelbegin \DIFdel{in }\DIFdelend \DIFaddbegin \DIFadd{on }\DIFaddend the surface buoyancy of the Arctic Ocean.

Figure \ref{fig:fig2}a and b show two snapshots of the buoyancy flux at the ocean surface. A negative buoyancy flux results from a loss of heat to the atmosphere sea ice refreezing and brine rejection, while a positive buoyancy flux is associated with sea ice melting\DIFdelbegin \DIFdel{and freshening}\DIFdelend \DIFaddbegin \DIFadd{, freshening, }\DIFaddend and warming of the Arctic surface. Inside the \DIFdelbegin \DIFdel{ice covered }\DIFdelend \DIFaddbegin \DIFadd{ice-covered }\DIFaddend region, the fluxes are spatially homogeneous, except for the marginal ice zone and within the leads, where the buoyancy fluxes are larger. The average buoyancy flux across the ice covered Arctic (15\% - 100\% ice concentration) within the leads is $-2.2\times 10^{-8} m^2/s^3$ and $7.2\times 10^{-8} m^2/s^3$ for each date, compared to the $-2.4\times 10^{-8} m^2/s^3$ and $5.7\times 10^{-8} m^2/s^3$ excluding the leads. In other words, on the 1st of January (winter), the total fluxes within the leads and those outside of the leads are comparable. However, on the 25th of June (summer), the average fluxes in the leads can be up to 30\% larger than those over the ice cover region excluding leads.

\begin{figure}
    \centering
    \includegraphics[width=1\textwidth]{./Figures/Figure_2_buoyancy_flux_maps_vertical_V2_15p.pdf}
    \caption{Maps of buoyancy flux and transects for the 1st of January and 25th of June 2014. a) Map of buoyancy flux on the 1st of January 2014. 
    Transect on the 1st of January 2014 of c) ice thickness, e) downward heat flux, g) freshwater flux (positive upward), and i) buoyancy flux (heat flux minus freshwater flux). b) Map of buoyancy flux on the 25th of June 2014. Transect on the 25th of June 2014 of d) ice thickness, f) downward heat flux, h) freshwater flux (positive upward), and j) buoyancy flux. Note that positive buoyancy fluxes correspond to buoyancy loss, while negative buoyancy fluxes correspond to buoyancy gain by the ocean. The transect location and orientation changes between the winter and summer dates. Magenta contour in panels a and b \DIFdelbeginFL \DIFdelFL{correspond }\DIFdelendFL \DIFaddbeginFL \DIFaddFL{corresponds }\DIFaddendFL to the 15\% concentration contour of sea ice. Vertical blue bars \DIFdelbeginFL \DIFdelFL{shows }\DIFdelendFL \DIFaddbeginFL \DIFaddFL{show }\DIFaddendFL the mask of leads in the transect.}
    \label{fig:fig2}
\end{figure}

Examining a transect in the Canadian Basin on the 1st of January 2014, we detect two leads (vertical blue bars) with a significant decrease in sea ice thickness of $\sim 25cm$, suggesting an opening of the sea ice cover (Fig.\ref{fig:fig2} c).
This transect exhibits small heat flux at the ocean surface and within the identified leads (Fig.\ref{fig:fig2} e) \DIFdelbegin \DIFdel{, }\DIFdelend because the ocean surface is at freezing point and the heat exchanged with the atmosphere is used to form sea ice rather than cooling the ocean.
The freshwater flux (Fig.\ref{fig:fig2}g) exhibits two prominent positive peaks located within the identified leads. Both peaks are twice the magnitude of the fluxes outside the leads, highlighting enhanced sea ice growth and brine rejection within the leads. 
The total buoyancy flux is negative for this snapshot, with the largest buoyancy loss \DIFdelbegin \DIFdel{occurs }\DIFdelend \DIFaddbegin \DIFadd{occurring }\DIFaddend within the leads (Fig.\ref{fig:fig2} i).
On the 25th of June (Fig.\ref{fig:fig2}d), we detect several leads, but only some show a decrease of the sea ice thickness of more than $1m$ (Fig.\ref{fig:fig2} h). These specific leads also experience an increase in heat fluxes, reaching up to $\sim 25 W/m^2$ (Fig.\ref{fig:fig2}f) and a local decrease of the freshwater flux, indicating a freshening of the surface due to sea ice melt. This change is captured by the net positive buoyancy flux, decreasing the density at the surface (Fig.\ref{fig:fig2}j).
Of all the identified leads, some exhibit significant fluxes at the ocean surface, while others have little to no imprint in the thickness and fluxes. This may be explained by the timing of the sea ice break up, as subsequent snapshots show an increase of fluxes at the location of these identified leads (not shown). Alternatively, previous studies have shown that leads generated through shear have a limited buoyancy exchange with the atmosphere, but may still enhance fluxes in the ocean due to isopycnal upwelling or downwelling forced by the atmosphere-ocean stress within the leads \citep{Bourgault_flux_lkf_2020}.

Similar to observations, we identify fluxes that are larger within the leads than outside of them for these summer and winter days. In the next section, we explore the spatially averaged contribution of leads to the surface \DIFdelbegin \DIFdel{the }\DIFdelend buoyancy fluxes and the integrated influence on the water mass transformation across the Arctic Basin during 2014. 

\section{Integrated contribution of leads to the Arctic buoyancy}\label{sec3}

%DIF < The ice covered Arctic experiences two main dynamical sea ice regimes, the marginal ice zone (MIZ; 15-80\% sea ice concentration) and the ice pack (80-100\% sea ice concentration). Thus, 
The following analysis is conducted separately for four distinct masks representing regions with leads and regions excluding leads for the MIZ and the ice pack. On average over 2014, these masks account for 6\% for the MIZ, 72\% for the ice pack, 6\% for the leads in the MIZ, and 16\% for the leads in the ice pack. Thus the part covered by leads is $\sim 22\%$ of the total sea ice covered Arctic, consistent with previous model assessments \citep{Wang_leads_2016}. 
The extent of each mask exhibits a seasonal cycle (Fig. \ref{fig:fig3}a). From winter (December, January, and February) to  summer (June, July, and August), the ice pack extent decreases from approximately $80\%$ to $60\%$, while the extent of the MIZ increases from $2\%$ to approximately $13\%$. 
On one hand, the percentage of the ice pack leads extent peaks in winter at 16\% and decreases to a minimum extent of 13\% in summer. On the other hand, the extent of the MIZ leads covers \DIFdelbegin \DIFdel{2\% }\DIFdelend \DIFaddbegin \DIFadd{$2\%$ }\DIFaddend in winter and around \DIFdelbegin \DIFdel{13\% }\DIFdelend \DIFaddbegin \DIFadd{$13\%$ }\DIFaddend in summer, approximately the same as the MIZ extent. 

The contribution of the total buoyancy flux occurring within the leads and ice covered regions \DIFdelbegin \DIFdel{are }\DIFdelend \DIFaddbegin \DIFadd{is }\DIFaddend quantified by weighting the buoyancy flux of each of the sea ice cover masks by the total ice covered area (Fig. \ref{fig:fig3}b; note that the size of each mask evolves in time). The mean buoyancy flux under the ice pack in winter (December-February) and summer (June-August) is $-1.8 \times10^{-8} m^2/s^3$  and $2.2 \times10^{-8} m^2/s^3$, respectively. The buoyancy flux underneath the ice pack changes \DIFdelbegin \DIFdel{sign }\DIFdelend \DIFaddbegin \DIFadd{signs }\DIFaddend during the seasonal cycle, suggesting a shift from ice formation in winter to ice melting in summer. Note that in the ice pack, the haline component of the buoyancy flux determines the buoyancy flux, while the heat contribution is negligible (See Fig.\ref{fig:S_fig2}). Meanwhile, the mean buoyancy flux in the MIZ is negligible in winter and increases to $0.9 \times10^{-8}m^2/s^3$ in summer. The MIZ has predominantly a positive buoyancy flux suggesting preferential melting throughout the year (except for autumn), and the haline component also determines the buoyancy flux (See Fig.\ref{fig:S_fig2}). In both\DIFaddbegin \DIFadd{, }\DIFaddend the MIZ leads and the ice pack leads, the buoyancy flux \DIFdelbegin \DIFdel{follow }\DIFdelend \DIFaddbegin \DIFadd{follows }\DIFaddend the same seasonality of the fluxes outside of the leads. Winter values are $-0.3 \times10^{-8}m^2/s^3$ and $-0.01 \times10^{-8}m^2/s^3$, while summer values reach $0.6 \times10^{-8}m^2/s^3$ and $0.9 \times10^{-8}m^2/s^3$ for the contribution of the leads in the ice pack and the MIZ, respectively.
Overall, leads explain between 20\% to 45\% of the ice covered buoyancy fluxes (red dashed line in Fig \ref{fig:fig3}b), more than the area extent of all leads that ranges from 18\% in winter to 26\% in summer. Therefore, fluxes are larger within the leads relative to their area extent. 

%DIF <  The area weighted buoyancy fluxes by each region (Fig. \ref{fig:fig3}c) highlights the time periods in which fluxes are distinct within leads and those underneath the surrounding ice cover.
%DIF < During winter, the area weighted buoyancy flux of the MIZ leads is negative (-0.8 $\times10^{-8} m^2/s^3$), compared to a positive one outside the leads (1.5 $\times10^{-8} m^2/s^3$). Thus, in winter leads are re-freezing ice, while underneath the ice-covered MIZ ice is melting.  During summer, the area weighted buoyancy flux within the ice pack leads is $4.4 \times10^{-8} m^2/s^3$, approximately 20\% more than the buoyancy flux underneath the ice pack ($3.6 \times10^{-8} m^2/s^3$), indicating enhanced melting within the ice pack leads. 
\DIFdelbegin %DIFDELCMD < 

%DIFDELCMD < %%%
\DIFdelend \begin{figure}
    \centering
    \includegraphics[width=1\textwidth]{./Figures/Figure_3_buoyancy_flux_ts_only_total.pdf}
    \caption{\DIFdelbeginFL \DIFdelFL{Time-series }\DIFdelendFL \DIFaddbeginFL \DIFaddFL{Time series }\DIFaddendFL during 2014 of a) the percentage of ice cover extent and b) the total buoyancy fluxes for the total ice covered, the ice pack, the MIZ, and the leads in the ice pack and \DIFdelbeginFL \DIFdelFL{in }\DIFdelendFL the MIZ. The MIZ is defined between ice concentration of 15-80\%, the ice pack is defined where ice concentration $>80\%$, and the total ice covered is defined as the sum of the MIZ and ice pack.}
    \label{fig:fig3}
\end{figure}

The role and impact leads have on the surface water masses of the Arctic can be quantified by using the surface water mass transformation framework (WMT; see methods section; \citealp{Walin_WMT_1982}). This approach provides information on the transport of water through the surface density contours due to surface buoyancy forcing that results in the formation of lighter or denser water at the surface. Thus, the WMT relates the surface buoyancy forcing and the surface density field to the properties of the ocean interior. 
Figure \ref{fig:fig4}a shows the averaged WMT during 2014 for the total ice covered Arctic, the ice pack, the MIZ, and the parts of these regions covered by leads. The definitions of the water masses depicted in figure \ref{fig:fig4}a follow roughly those proposed by \citealt{Pemberton_WMT_2015} and \citealt{Lansard_water_mass_2012}.
The yearly mean WMT in the ice pack is characterized by a broad region of strong positive transformation (losing buoyancy) corresponding approximately to the halocline polar surface water ($PSW_{H}$; $24.4 kg/m^3 < \rho < 27.4 kg/m^3$) reaching a maximum of $\sim 1.8 Sv$, consistent with previous estimates of the surface Arctic WMT \citep{Pemberton_WMT_2015}.
The yearly mean WMT in the ice pack leads also features a positive peak within the $PSW_{H}$ density class reaching up to $0.4 Sv$. On average during 2014, the WMT in the ice pack leads is $\sim 20\%$ of the WMT in the ice pack (without leads). The mean MIZ WMT is mostly negative within the density range of $20 kg/m^3 < \rho < 27.5 kg/m^3$ (gaining buoyancy), reaching a minimum of $-0.4 Sv$ around $\rho = 27 kg/m^3$. The MIZ leads WMT is also negative within the same range of densities, with a minimum \DIFdelbegin \DIFdel{also }\DIFdelend of $-0.4 Sv$. Thus, the MIZ leads WMT is comparable to those in the MIZ excluding leads. 
Hovm\"oller diagrams of the daily WMT estimates over 2014 for the four masks are depicted in figure \ref{fig:fig4}.
The mean WMT underneath the ice pack varies from 1.1 Sv in winter to -1.0 Sv in summer, while the WMT in the leads in the ice pack is 0.2 Sv in winter and -0.26 Sv in summer. In other words, the WMT in the ice pack leads ranges from 20 to 25\% of the WMT occurring in the pack (Fig. \ref{fig:fig4}d). It is interesting to note that the patterns across the density classes of the WMT for the leads and the surrounding regions excluding leads are similar (figures \ref{fig:fig4}b and c). This suggests that there is \DIFdelbegin \DIFdel{not }\DIFdelend \DIFaddbegin \DIFadd{no }\DIFaddend specific water mass that is transformed preferentially with the leads. Rather, our results \DIFdelbegin \DIFdel{suggests }\DIFdelend \DIFaddbegin \DIFadd{suggest }\DIFaddend that the water masses transformed within the leads and outside of them have similar properties. 
The WMT in MIZ leads and those in the MIZ WMT have comparable magnitudes over the year (Fig. \ref{fig:fig4} g). Similarly to the leads in the ice pack, there is no evidence that leads in the MIZ could transform preferentially specific water masses compared to the remaining MIZ. 
Adding the contribution of the WMT in all the leads in the Arctic (those in the MIZ and the ice pack), we estimate $\sim 25\%$ of the WMT happens in leads, and this is comparable to the surface covered by these leads in the Arctic ($\sim 22\%$). 

\begin{figure}[!ht]
    \centering
    \includegraphics[width=1\textwidth]{./Figures/Figure_4_water_mass_transformation_hovmoller_V2_MIZ15_80p.pdf}
    \caption{Mean and daily Water Mass Transformation (WMT) over 2014. a) Mean water mass transformation over 2014 for the total ice covered, ice pack, MIZ, and leads in each region. Hovm\"oller diagram of the water mass transformation (Sv) for b) the ice pack, c) the ice pack leads\DIFaddbeginFL \DIFaddFL{, }\DIFaddendFL and d) the ratio of water mass transformation between the leads and the ice pack region.
    Panels e), f), and g) follow the same structure as panels b), c)\DIFaddbeginFL \DIFaddFL{, }\DIFaddendFL and d), but for the MIZ, MIZ leads, and \DIFaddbeginFL \DIFaddFL{the }\DIFaddendFL ratio between them.
    Shaded gray and white areas in panel a) represent approximately the water masses of the Arctic \DIFdelbeginFL \DIFdelFL{ocean }\DIFdelendFL \DIFaddbeginFL \DIFaddFL{Ocean }\DIFaddendFL according to \citealt{Pemberton_WMT_2015} and \citealt{Lansard_water_mass_2012}: the mixed layer polar surface water ($PSW_{ML};\ \rho < 24.4 kg/m^3$), the halocline polar surface water ($PSW_{H};\ 24.4 kg/m^3 < \rho < 27.4 kg/m^3$), the polar deep water ($PDW;\ 27.4 kg/m^3 < \rho < 27.8 kg/m^3$), and the Atlantic water ($AW;\ 27.8 kg/m^3 < \rho$).}
    \label{fig:fig4}
\end{figure}

\section{Conclusions}
%DIF < The lead identification algorithm implemented in SEDNA allows us to reproduce the observed spatial distribution of leads over most regions of the Arctic, with exception of the Beaufort Sea, where the simulation underestimates the lead frequency. Likely because of the different identification algorithm, time periods, and spatial and temporal scales used in the present study and that of \citealp{Willmes_leads_currents_2023}. Despite these regional differences, SEDNA can estimate the mean lead probability, showcasing its capacity to capture the dynamics of these features.

The present study assesses the impact of leads \DIFdelbegin \DIFdel{in }\DIFdelend \DIFaddbegin \DIFadd{on }\DIFaddend the ocean surface by diagnosing the buoyancy fluxes using a very high-resolution hindcast (SEDNA). 
SEDNA captures the intermittency of the leads and the lead properties. Leads cover $\sim 22\%$ of the sea-ice covered Arctic, and in the ice pack leads can contribute more buoyancy flux compared to the fluxes outside the leads. Despite these larger buoyancy fluxes, leads only account for approximately $\sim 25\%$ of the total surface water mass transformation in the ice covered Arctic, which is roughly equivalent to their surface coverage. Thus, contrary to previous hypotheses based largely on local and intermittent observations (e.g. \citet{Hutter_identification_2019, McPhee_LKF_2005}), our findings suggest that leads have a smaller contribution to the transformation of surface water masses in the Arctic Ocean. While their presence can induce localized changes in the density of surface waters due to brine rejection and ice melting, 
the transformation of water masses in the leads \DIFdelbegin \DIFdel{have }\DIFdelend \DIFaddbegin \DIFadd{has }\DIFaddend the same sign and magnitude as those underneath the surrounding ice-covered regions. Thus, no evidence was found that leads impact in a distinct way the Arctic \DIFdelbegin \DIFdel{ocean }\DIFdelend \DIFaddbegin \DIFadd{Ocean }\DIFaddend water masses in SEDNA compared to outside the leads in the ice covered ocean.

Here we solely focus on the surface forced WMT, but we acknowledge that leads could further impact properties at the surface and interior of the ocean by enhancing mixing within leads, inducing upwelling of isopycnals through Ekman pumping \citep{McPhee_LKF_2005}, modifying  entrainment in the mixed layer within the leads, providing a source of potential energy, and generating mixed layer instabilities due to isopycnal tilting in the vicinity of leads. The use of a continuous sea ice model that results in leads with high ice thickness and concentrations likely \DIFdelbegin \DIFdel{underestimate }\DIFdelend \DIFaddbegin \DIFadd{underestimates }\DIFaddend the exchanges between the ocean and the atmosphere. Therefore, future studies using \DIFdelbegin \DIFdel{floe resolving }\DIFdelend \DIFaddbegin \DIFadd{floe-resolving }\DIFaddend models that allow more ``realistic leads" with direct exchanges between the ocean \DIFdelbegin \DIFdel{an }\DIFdelend \DIFaddbegin \DIFadd{and }\DIFaddend the atmosphere may result in larger fluxes through the leads increasing their potential impact in the water mass transformation and ocean mixed layer.

Leads have an important local effect in the atmosphere and the sea-ice fluxes \citep{Marcq_flux_lead_2012, Lupkes_leads_2008, Wettlaufer_insulator_1997, Smith_polynyas_1990}, but a limited imprint on the ocean surface and interior properties. Our results suggest that the presence of leads is not the first order mechanism controlling the water mass transformation forced at the surface, thus leads have a minor impact \DIFdelbegin \DIFdel{in }\DIFdelend \DIFaddbegin \DIFadd{on }\DIFaddend the stratification and surface layer dynamics of the Arctic Ocean.
Therefore, resolving and/or parameterizing leads in climate models are likely to only marginally improve the misrepresentation of the stratification and water masses found in Arctic climate models \citep{Wang_ompi_2024,Ilicak_core_2016}. Finally, the probability of leads \DIFdelbegin \DIFdel{to occur }\DIFdelend \DIFaddbegin \DIFadd{occurring }\DIFaddend in the ice pack has been projected to increase as the Arctic sea ice transitions towards thinner, younger, and more mobile sea ice \citep{Boutin_rheology_2023, IPCC_arctic_2022}. Our analysis is only performed over one year and the future impact of more abundant leads is not represented in our simulation. Thus, dedicated studies are required to better estimate the surface buoyancy forcing within leads in the context of an evolving Arctic Ocean.

\section{Methods}
\subsection*{Model description}

Here we employ the NEMO numerical platform (release 4.0.5; \citealp{NEMO_man}) to build a 1/60$^\circ$ pan-Arctic ocean-sea ice configuration named SEDNA. The ocean core (NEMO-OPA) solves the primitive equations using finite differences on an Arakawa C-grid mesh, and it is coupled to the SI3 sea-ice component (NEMO-SI3; \citealp{SI3_man}). The pan-Arctic domain covers the Bering Strait on the Pacific side and extends southward to 56$^\circ$N in the Subpolar North Atlantic and to the Baltic Sea entrance. The ocean model incorporates a linear free surface, utilizes a 3rd order flux form scheme for momentum advection, and employs a 4th order Flux Corrected Transport (FCT) for tracer advection. Additionally, the sea-ice representation uses an Elasto-Visco-Plastic (EVP) rheology with a 5-categories sea ice thickness distribution and a landfast ice parameterization to better simulate sea-ice behavior above shelves.

The NCAR bulk formula from \citet{Large_global_2009} is used to calculate air-sea fluxes based on the hourly ERA5 atmospheric data for near-surface variables \citep{Hersbach_ERA5_2020}. The three lateral boundaries are constrained with daily GLORYS12V1 Re-analysis data \citep{GLORYS12}. Monthly freshwater fluxes from river discharges are combined with Greenland land ice melt data \citep{HuEA_rivers_2019}. The simulation extends over seven years (2009 to 2015) and starts from rest, utilizing initial conditions based on the World Ocean Atlas 2009 temperature/salinity and mean January 2009 ice state from PIOMAS re-analysis \citep{Locarnini_data_2010, PIOMAS}. A weak restoring toward the World Ocean Atlas climatological sea surface salinity is applied, but only over the ice-free regions. Preliminary assessment against available observations and re-analyses data shows that the simulation reasonably captures key ocean and sea-ice characteristics.

\subsection*{Identification of leads}

The leads from SEDNA are identified by using the algorithm proposed by \citet{Hutter_identification_2019}. This method uses the total deformation rate defined as: 
\begin{equation}
    T_d = \left[\left(\underbrace{\frac{\partial u}{\partial x} +\frac{\partial v}{\partial y}}_{Divergence}\right)^2 + \underbrace{\left(\frac{\partial u}{\partial x} - \frac{\partial v}{\partial y}\right)^2+ \left(\frac{\partial u}{\partial y} + \frac{\partial v}{\partial x}\right)^2}_{Shear}\right]^{1/2},
\end{equation}
where $u$ and $v$ are the sea ice velocities. The deformation rate varies spatially depending on the background deformation and the sea ice properties \citep{Reiser_LKF_2020}, where the largest deformation rates occur along the boundaries of ice floes where the leads are found. Therefore, to identify the leads, the local maxima is identified using a difference of Gaussian filter (DoG filter) of the deformation rate field. After this, a binary map is created with the mask of identified lead. Note that after this step, \citet{Hutter_identification_2019} reduce the leads width to a line to then track the leads over time. Since our study \DIFdelbegin \DIFdel{focus }\DIFdelend \DIFaddbegin \DIFadd{focuses }\DIFaddend on the ocean-atmosphere interactions, we only use the identification algorithm to produce a mask of leads during 2014 for the Arctic, thus we retain the original lead width computed from the deformation rate. Finally, the lead mask is then used to quantify and contrast the impact of leads compared to the ice covered Arctic, the ice-pack (ice concentration $> 80\%$)\DIFaddbegin \DIFadd{, }\DIFaddend and the MIZ (0 - 80\% ice concentration).

\subsection*{Buoyancy flux and water mass transformation}

Heat fluxes (shortwave radiation, longwave radiation, latent heat, and sensible heat) and freshwater fluxes (evaporation, precipitation, runoff, and sea ice melt/growth) at the ocean surface result in a buoyancy forcing capable \DIFdelbegin \DIFdel{to change }\DIFdelend \DIFaddbegin \DIFadd{of changing }\DIFaddend the density of the ocean surface. Here, the buoyancy flux is computed as:

\begin{equation}
    B_f = \underbrace{\frac{g \alpha Q_{net}}{\rho_0 C_p}}_{Net\ surface\ heat\ flux} - \underbrace{\frac{g \beta F_{net} SSS }{\rho_0}}_{Net\ surface\ freshwater\ flux},
\end{equation}

where $Q_{net}$ the net heat flux at the sea surface, $\alpha$ the thermal expansion coefficient, $SSS$ the sea surface salinity, $F_{net}$ the net \DIFdelbegin \DIFdel{fresh water }\DIFdelend \DIFaddbegin \DIFadd{freshwater }\DIFaddend flux that includes evaporation, precipitation, runoff, and freezing minus melting flux, \DIFdelbegin \DIFdel{and }\DIFdelend $\beta$ the haline contraction coefficient, \DIFdelbegin \DIFdel{and }\DIFdelend $g = 9.81 m/s^2$\DIFaddbegin \DIFadd{, }\DIFaddend and the ocean density $\rho_0 = 1026 kg/m^3$. 
Note that using this sign convention, positive values in the buoyancy forcing decrease the surface density making the ocean surface more buoyant (stable) and negative values increase the surface density. The area-weighted buoyancy ($ \left<\right>$) fluxes are computed following:

\begin{equation}
    \left<B_f (t)\right>  = \frac{\sum_x \sum_y \left( B_f(t,x,y) * Area(x,y) * Mask_{_R}(t,x,y)\right)}{\sum_x \sum_y * Area(x,y) * Mask(t,x,y)},
\end{equation}
where the $Area$ is \DIFdelbegin \DIFdel{be }\DIFdelend the area of the grid cells, and the $Mask_{_R}$ can be the mask of MIZ, the Pack, the MIZ leads, the ice pack leads, and the ice cover that \DIFdelbegin \DIFdel{vary }\DIFdelend \DIFaddbegin \DIFadd{varies }\DIFaddend in space and time. Note that the divisor $Mask$ corresponds to the full ice covered area (including leads).

\citet{Nurser_WMT_1999} and  \citet{Walin_WMT_1982} proposed the water mass transformation framework that quantifies the rate at which a given water mass is made lighter (negative transformation rate) or denser (positive) due to surface fluxes. The surface water mass transformation is computed as:

\begin{equation}
    \Omega(\sigma_k) = - \frac{1}{\sigma_{k+1}-\sigma_k}\iint_AB_fdA,
\end{equation}

%DIF <  \begin{equation}
%DIF <      \Omega(\sigma_k) = \underbrace{- \frac{1}{\sigma_{k+1}-\sigma_k}\iint_A\left(\frac{\alpha Q_{net}}{\rho_0 C_p}\right)dA}_{Net\ surface\ heat\ flux} + \underbrace{\frac{1}{\sigma_{k+1}-\sigma_k}\iint_A\left(\frac{\beta F_{net} SSS }{\rho_0}\right)dA}_{Net\ surface\ freshwater\ flux},
%DIF <  \end{equation}
\DIFdelbegin %DIFDELCMD < 

%DIFDELCMD < %%%
\DIFdelend where $\sigma$ is potential density and the indices $k$ represents the density bin number. The computation was performed with 172 density bins of $0.1 kg/m^3$ within the density range of $15.15 kg/m^3$ - $32.25 kg/m^3$. The transformation rate $\Omega$ is spatially decomposed into different sea-ice masks such as the ice covered, the pack ice, the MIZ, and the leads in each of these regions. 
%DIF < By separately diagnosing the density transformation that occurs in the leads and underneath the surrounding floes, we quantify the relative importance of the leads for each of these regions. Additionally, this framework accounts for the area changes of the ice cover during the seasonal cycle, and excludes open ocean regions.

\backmatter

\bmhead{Open Source}

The SEDNA model configuration are described and publicly available at \url{https://zenodo.org/records/7656924}.
All analyses and figures in this manuscript are reproducible via Jupyter notebooks and instructions can be found in the GitHub repository \texttt{Leads\_Arctic\_ocean} (\url{https://github.com/josuemtzmo/Leads_Arctic_ocean}).

\bmhead{Acknowledgements}

We acknowledge funding from the ANR ImMEDIAT project (ANR-18-CE01-0010) and the MEDLEY project funded by the program JPI Ocean/JPI Climate (ANR-19-JPOC-0001) project. 
The idealized simulations and their analysis were performed using the HPC facilities DATARMOR of \DIFdelbegin \DIFdel{’}\DIFdelend \DIFaddbegin \DIFadd{`}\DIFaddend Pôle de Calcul Intensif pour la Mer\DIFdelbegin \DIFdel{’ }\DIFdelend \DIFaddbegin \DIFadd{' }\DIFaddend at Ifremer, Brest, France. We also acknowledge PRACE for awarding us access to Joliot-Curie at GENCI\makeatletter@CEA, France, where the SEDNA simulation has been performed.

%%===========================================================================================%%
%% If you are submitting to one of the Nature Portfolio journals, using the eJP submission   %%
%% system, please include the references within the manuscript file itself. You may do this  %%
%% by copying the reference list from your .bbl file, paste it into the main manuscript .tex %%
%% file, and delete the associated \verb+\bibliography+ commands.                            %%
%%===========================================================================================%%

%% if required, the content of .bbl file can be included here once bbl is generated
%%\input sn-article.bbl


\bibliography{biblio}


\pagebreak

\begin{appendices}

\renewcommand{\thefigure}{S\arabic{figure}}

\begin{figure}
    \centering
    \includegraphics[width=1\textwidth]{./Figures/Figure_S1_comparison_LKF_observations_KE_bathy.pdf}
    \caption{Comparison of winter lead frequency between a) Thermal-Infrared Satellite Imagery \citep{Willmes_leads_currents_2023} and b) SEDNA between 2011 and 2015. Panel c) shows the SEDNA mean kinetic energy between 2011 and 2015.}
    \label{fig:S_fig1}
\end{figure}

\begin{figure}
    \centering
    \includegraphics[width=1\textwidth]{./Figures/Figure_S2_buoyancy_flux_components.pdf}
    \caption{Time-series during 2014 of a) the freshwater component of the buoyancy flux and b) the heat component of the buoyancy flux for the ice pack, the MIZ, and the leads in the ice pack and in the MIZ. The MIZ is defined between ice concentration of 15-80\%, and the ice pack is defined where ice concentration $>80\%$.}
    \label{fig:S_fig2}
\end{figure}

\end{appendices}

\end{document}
